\section{Study 3: the automatic mode at work (tracetrim 0.4)}
\label{sec:study3}

Studies~1 and~2 learned a projection offline and applied it fixed. Version 0.4 of the library added what this study
measures. In \emph{automatic mode} the library learns while the agent works: runs are recorded in a work directory,
a run counts for learning once the agent's final answer is known (through \texttt{tracetrim\_done} behind the MCP
proxy, or directly in a function-calling loop), and trimming starts once enough runs are recorded. Each
\texttt{expand} request is kept as evidence, with two optional corrections: \emph{promotion}, which keeps a field
that the evidence shows to be needed, and \emph{sampling}, which passes a share of the results that would be
trimmed whole. A tool whose trimmed results are expanded often enough stops being trimmed (\emph{switch-off}).
Version 0.4 also ships a \emph{skill}, a text for the agent: its base level explains \texttt{tracetrim\_done} and
\texttt{expand}; an optional \emph{strategy} level teaches how to organise tool calls (independent calls together,
grouped inputs, when to stop) and prescribes no parameter value. Study~3 changes nothing in Studies~1 and~2.

\paragraph{Design.} Two protocols, each frozen before its first paid run. \texttt{protocol-v04.md} (frozen on
\SthreeFrozenVfour) covers E1, on the strategy level of the skill, and E2, on adaptive learning:
\SthreeRunsEOneTwo{} runs of \texttt{\SthreeModelA}, the model of Studies~1 and~2, at default parameters.
\texttt{protocol-v04-e3.md} (frozen on \SthreeFrozenEthree) covers E3, the strategy on a second model,
\texttt{\SthreeModelB}: \SthreeRunsEThree{} runs after \SthreeRunsEThreePilot{} pilot runs that enter no
result. All runs used the library's function-calling adapter in the study's own loop; the MCP proxy was not
exercised (that the two adapters take the same decisions is covered by a test, not by these protocols). The
protocols list \SthreeDeviationsVfour{} and \SthreeDeviationsEthree{} deviations respectively, dated and explained,
each decided before the measured runs it concerns; two of them reduced E2 on the real-API case to two arms of one
sequence each. The protocols number their hypotheses H1--H11 independently of Studies~1 and~2; in this section
H-numbers refer to them. As before, a difference inside the other arm's range is not read as a difference, and no
significance test is claimed. Table~\ref{tab:s3hyp} gives every outcome as frozen.

\begin{table}[t]
\centering\small
\resizebox{\textwidth}{!}{\begin{tabular}{llp{6.2cm}lp{5.2cm}}
\toprule
& Exp. & Hypothesis, as frozen & Outcome & Detail \\
\midrule
P4-H1 & E1 & With the strategy, gates and coverage are no worse & failed & 4 of 5 comparisons hold; editorial plan: ST passed the gates in 4/5 runs, T in 5/5 \\
P4-H2 & E1 & Billing support: median turns with the strategy below the minimum without & failed & median 4 against minimum 3 \\
P4-H3 & E1 & Removed before freezing (it concerned arms that are not run) & -- & -- \\
P4-H4 & E1 & The strategy does not undo the saving of trimming (billed input) & failed & editorial plan holds; billing support fails \\
P4-H5 & E2 & A field recovered with \texttt{expand} and written is kept from the next run & holds & 16 of 16 recoveries, 10 sequences \\
P4-H6 & E2 & A reworded field needs promotion to come back & not testable & reworded in 2 sequences, all with promotion on \\
P4-H7 & E2 & Sampling reaches a field the agent does not ask for & failed & runs with a field removed and no \texttt{expand}: 1 with sampling, 1 without; requirement met in both \\
P4-H8 & E3 & Median turns with the strategy below the minimum without & failed & holds in 2 of 4 cases \\
P4-H9 & E3 & Median share of turns with several calls above the maximum without & failed & holds in 0 of 4 cases \\
P4-H10 & E3 & Gates and coverage are no worse with the strategy & holds & gates, billing support: 8/8 with, 8/8 without; gates, CI triage: 6/6 with, 6/6 without \\
P4-H11 & E3 & The strategy does not override the values the prompts prescribe & holds & 9 of 9 checks \\
\bottomrule
\end{tabular}
}
\caption{Study 3: the hypotheses of \texttt{protocol-v04.md} (E1, E2) and \texttt{protocol-v04-e3.md} (E3), stated
before the runs and numbered independently of Studies~1 and~2, with their outcomes as frozen. A hypothesis whose
precondition did not occur is \emph{not testable}. ST and T: arms of E1 with the strategy and trimming, and with
trimming alone.}
\label{tab:s3hyp}
\end{table}

\subsection{E1: the strategy level of the skill}
\label{sec:s3e1}

Four arms: R (neither strategy text nor trimming), S (strategy text), T (trimming with \texttt{expand}) and ST
(both), \SthreeRunsPerArmEOne{} runs each, on the editorial plan, the content-gap analysis and billing support,
with the prompts of Studies~1 and~2; \SthreeRunsEOne{} runs in all. Trimming is a fixed configuration learned from
the \PhaseARunsi{} Phase A runs of the case with the library's default rules; on the content-gap case these rules
trim no tool, so it has arms R and S only.

\paragraph{The strategy changed nothing measurable on this model.} Against the arm without it, the median of S
stayed inside R's range for turns, tool calls, billed input tokens and the share of turns with several calls in
\SthreeEOneSinsideR{} of \SthreeEOneSinsideRof{} comparisons (\SthreeEOneCases{} cases, \SthreeEOneMeasures{} measures;
Table~\ref{tab:s3e1}).
There was little to organise: without the strategy, \SthreeEOneSeveralRi, \SthreeEOneSeveralRii{} and
\SthreeEOneSeveralRiv{} of the turns that made a call made several (editorial plan, content gaps, billing support;
\SthreeEOneSeveralSi, \SthreeEOneSeveralSii{} and \SthreeEOneSeveralSiv{} with it). The hypotheses on the strategy
failed as frozen. H1 failed in \SthreeHOneFails{} comparison of \SthreeHOneOf: in the editorial plan ST passed the gates in
\SthreeHOneGatesST{} runs against \SthreeHOneGatesT{} for T. H2 failed: on billing support the median turns of S were \SthreeHTwoS, not below R's minimum of \SthreeHTwoR. H4 failed on billing support, where the median billed input of ST (\SthreeEOneBilledSTiv{} tokens) was
above T's maximum (\SthreeEOneBilledTmaxiv); in the editorial plan it was below T's minimum (\SthreeEOneBilledSTi{}
against \SthreeEOneBilledTmini). With the strategy, the values the prompts prescribe (result limits, search depth,
keywords per call) were \SthreeEOnePrescribedKept{} in all \SthreeEOnePrescribedRuns{} runs of the two cases whose
prompts prescribe them. Runs that asked a question instead of doing the task: \SthreeEOneAsked{} of the
\SthreeEOneStrategyRuns{} with the strategy (in arm ST), \SthreeEOneAskedOther{} of those without it.

\paragraph{Trimming, measured as a side arm.} E1 was not about trimming and states no hypothesis on it, but its
arm T gives a live measure of the fixed configuration against R. The tool outputs as passed to the model were
\SthreeEOneTrimTokensi{} smaller than the same outputs unfiltered in the editorial plan, and \SthreeEOneTrimTokensiv{}
smaller in billing support (medians of the runs); the median cost of a run, at the model's list prices, was
\SthreeEOneTrimCosti{} and \SthreeEOneTrimCostiv{} lower than in R; the gates passed in \SthreeEOneTrimGatesi{} and
\SthreeEOneTrimGatesiv{} runs. Costs are reported only as changes relative to R.

\begin{table}[t]
\centering\small
\resizebox{\textwidth}{!}{\begin{tabular}{llrrlllrrr}
\toprule
Case & Arm & Runs & Gates & Turns & Tool calls & Turns with several calls & Billed input & Tool-output tokens & Cost of a run \\
\midrule
Editorial plan & R & 5 & 4/5 & 9 [8--10] & 22 [19--24] & 75\% [67--88] & 704{,}930 & -- & -- \\
 & S & 5 & 5/5 & 10 [6--11] & 20 [18--26] & 80\% [56--100] & 763{,}969 & -- & +2\% \\
 & T & 5 & 5/5 & 9 [8--11] & 24 [17--26] & 88\% [60--100] & 476{,}266 & $-$44\% & $-$30\% \\
 & ST & 5 & 4/5 & 9 [1--12] & 23 [0--27] & 81\% [64--89] & 442{,}252 & $-$44\% & $-$35\% \\
\addlinespace
Content gaps & R & 5 & 3/5 & 5 [3--6] & 16 [14--20] & 67\% [60--100] & 236{,}839 & -- & -- \\
 & S & 5 & 4/5 & 4 [3--6] & 16 [13--18] & 100\% [50--100] & 191{,}175 & -- & $-$7\% \\
\addlinespace
Billing support$^\dagger$ & R & 5 & 5/5 & 4 [3--4] & 11 [11--13] & 67\% [67--100] & 43{,}275 & -- & -- \\
 & S & 5 & 5/5 & 4 [3--4] & 12 [10--13] & 67\% [67--100] & 45{,}359 & -- & $-$1\% \\
 & T & 5 & 5/5 & 3 [3--4] & 10 [10--11] & 100\% [67--100] & 17{,}340 & $-$44\% & $-$24\% \\
 & ST & 5 & 5/5 & 4 [3--4] & 11 [10--12] & 67\% [67--100] & 33{,}613 & $-$43\% & $-$20\% \\
\bottomrule
\end{tabular}
}
\caption{E1, \texttt{\SthreeModelA}, \SthreeRunsPerArmEOne{} runs per arm. Turns and tool calls: median [min--max];
turns with several calls: share of the turns with a call that made more than one, per run, median [min--max] (the
text gives the share over all the turns of an arm). Billed input: median tokens of a run, as billed. Tool-output
tokens: the outputs as passed on, against the same outputs unfiltered. Cost of a run: median, relative to R, at
the model's list prices. $^\dagger$Simulated tools.}
\label{tab:s3e1}
\end{table}

\subsection{E2: adaptive learning}
\label{sec:s3e2}

Each case was given a variant of its task that needs two fields its projection removes: in billing support, the
link to an invoice and how it is collected; in the editorial plan, a video's channel and its views. A sequence is
\SthreeETwoSequenceLength{} consecutive runs in one work directory, seeded with the \PhaseARunsi{} Phase A runs
of the case, with automatic mode on and the base level of the skill. Billing support ran on a copy of its simulated
world in which those two invoice fields vary, with promotion off or on and sampling off or on, two sequences per
arm (\SthreeRunsETwoSim{} runs); the editorial plan ran the arms without promotion, base and sampling, one sequence
each (\SthreeRunsETwoReal{} runs).

\paragraph{Fields came back by themselves (H5 holds).} Each time the model recovered a needed field and wrote its
value, the field was in the trimmed outputs from the next run on, with no \texttt{expand}: \SthreeHFiveKept{} of
\SthreeHFiveRecoveries{} recovered and written fields were kept in every later run, in \SthreeETwoSequences{}
sequences; from the second run in \SthreeETwoFromSecond{} sequences and from the third in \SthreeETwoFromThird. The
requirement the variant adds was met in \SthreeETwoRequirementMet{} of \SthreeETwoRuns{} runs.

\paragraph{Promotion and sampling made no difference here.} Promotion exists for a field that is used without
being written; this model wrote the values it recovered, so ordinary learning was enough. The field expected to be
reworded was reworded in \SthreeHSixReworded{} sequences, both with promotion on, so H6 could not be tested.
Sampling only cost tokens: in the last run of each billing-support sequence, the median saving on tool-output tokens
was \SthreeETwoSavingSampling{} with sampling and \SthreeETwoSavingNoSampling{} without. The model never needed it:
H7 failed: in the runs where a needed field was removed and the model did not call \texttt{expand}
(\SthreeHSevenPerGroup{} with sampling, \SthreeHSevenPerGroup{} without), the requirement was met in both, so
not more often with sampling.

\paragraph{Another route to a removed field.} In those \SthreeETwoOtherTool{} billing-support runs the model did not
call \texttt{expand}: it called another tool, which retrieves a single invoice and is not trimmed, and read the
fields there. That is a limit of the design on that case (it left H7 little to test), and a behaviour to know: on an
API billed per call, that route has a price. In the editorial plan no other tool returns the two fields, and the
model used \texttt{expand} in both sequences.

\paragraph{One switch-off.} In the first run of the editorial plan's base sequence the model expanded
\SthreeSwitchExpanded{} of the \SthreeSwitchFiltered{} trimmed results of the trends tool, all in
\SthreeSwitchRunsWithExpand{} run, and under the rule then in force the tool stopped being trimmed. The saving of
that sequence was \SthreeSwitchSavingFirst{} in its first run and between \SthreeSwitchSavingLaterMin{} and
\SthreeSwitchSavingLaterMax{} later; in the sampling sequence, where no tool was switched off, it stayed between
\SthreeSwitchSamplingMin{} and \SthreeSwitchSamplingMax. An offline recalculation on the \SthreeETwoSequences{}
sequences, not in the protocol, replayed the rule with expansions required from \SthreeSwitchRuleA{} or
\SthreeSwitchRuleB{} distinct runs: \SthreeSwitchWithRules{} tool would have been switched off at that point (the
recalculation is exact up to the first run in which two rules decide differently; what the agent would have asked
afterwards is not in the traces). The library now requires expansions from \SthreeSwitchDefaultNow{} distinct runs
by default.

\paragraph{Requests that gave nothing back.} Of the \SthreeExpandRequests{} \texttt{expand} requests of E1 and E2,
\SthreeExpandNothing{} returned nothing. \SthreeExpandUnknownId{} named a result that did not exist (an identifier
made up or cut short), \SthreeExpandRows{} of them asking for more rows than the tool had returned: the agent wanted more rows, not hidden
fields, which \texttt{expand} cannot give.
The other \SthreeExpandIndex{} wrote a list index where paths collapse indices; the library now accepts such a path,
and answers an unknown identifier by saying how to copy a valid one.

\subsection{E3: the strategy on a second model}
\label{sec:s3e3}

E1 left open what the strategy does for a model that issues calls one at a time. E3 asked it of
\texttt{\SthreeModelB}, with arms R and S and no trimming, on billing support and CI triage, whose simulated
servers were live, and on the editorial plan and the content-gap analysis, whose search-data provider was replaced
by recorded outputs: a call got the recorded output of the same call when there was one and otherwise the output of
the recorded call of the same tool with the most similar arguments. The provider was called in \SthreeProviderCallsEThree{} run.
On those two cases only the organisation of calls is evaluated, not the answers. Runs per arm go from
\SthreeEThreeRunsPerArmMin{} to \SthreeEThreeRunsPerArmMax{} (Table~\ref{tab:s3e3}).

\begin{table}[t]
\centering\small
\resizebox{\textwidth}{!}{\begin{tabular}{llrlllrr}
\toprule
Case & Arm & Runs & Turns & Tool calls & Turns with several calls & Billed input & Cost of a run \\
\midrule
Billing support$^\dagger$ & without & 8 & 3.5 [3--8] & 11.5 [10--14] & 88\% [43--100] & 25{,}316 & -- \\
 & strategy & 8 & 3 [3--3] & 10 [10--12] & 100\% [100--100] & 17{,}196 & $-$20\% \\
\addlinespace
CI triage$^\dagger$ & without & 6 & 17 [5--19] & 16 [13--18] & 0\% [0--100] & 865{,}862 & -- \\
 & strategy & 6 & 6 [5--8] & 16.5 [15--20] & 90\% [57--100] & 391{,}018 & $-$15\% \\
\addlinespace
Content gaps$^\ddagger$ & without & 3 & 6 [6--7] & 15 [14--16] & 100\% [100--100] & 169{,}493 & -- \\
 & strategy & 3 & 5 [4--7] & 14 [13--19] & 100\% [100--100] & 131{,}489 & +1\% \\
\addlinespace
Editorial plan$^\ddagger$ & without & 2 & 14 [13--15] & 23 [21--25] & 47\% [36--58] & 1{,}023{,}210 & -- \\
 & strategy & 2 & 10.5 [10--11] & 18 [17--19] & 53\% [50--56] & 658{,}204 & $-$2\% \\
\bottomrule
\end{tabular}
}
\caption{E3, \texttt{\SthreeModelB}, without and with the strategy text, no trimming. Columns as in
Table~\ref{tab:s3e1}; cost relative to the arm without the strategy, at the model's list prices. Billed tokens are
not comparable with E1's (another tokenizer, price list and caching). $^\dagger$Simulated tools, live.
$^\ddagger$Recorded outputs in place of the provider: calls and turns only.}
\label{tab:s3e3}
\end{table}

\paragraph{The hypotheses on turns and grouped calls failed as stated.} H8 (median turns with the strategy below
the fewest turns of any run without it) held in \SthreeHEightHolds{} cases of \SthreeHEightOf, the two with
recorded outputs; H9 (median share of turns with several calls above the highest share of any run without it) held
in \SthreeHNineHolds{} case: in each, a run without the strategy reached the same share or more. Why they failed
concerns the rule, not the runs, whose results stand as they are. Comparing the median of one arm with the extreme
of the other lets a single run decide: with \SthreeEThreeRunsCv{} runs per arm, as on CI triage, one unusual run without the strategy
is enough to fail the hypothesis, and there that is what happened. A protocol that wants to know whether an
arm is usually shorter should compare what is usual in both. The two simulated cases are therefore read apart.

\paragraph{CI triage, where this model issued one call per turn.} Without the strategy, \SthreeEThreeLongRunsv{}
runs of \SthreeEThreeRunsCv{} made every call in a turn of its own (\SthreeEThreeLongMinv{} to
\SthreeEThreeLongMaxv{} turns); the remaining one grouped its calls and took \SthreeEThreeTurnsRminv. With the strategy all
runs took \SthreeEThreeTurnsSminv{} to \SthreeEThreeTurnsSmaxv{} turns, with a similar number of calls. The median
went from \SthreeEThreeTurnsRv{} turns to \SthreeEThreeTurnsSv, billed input tokens fell by
\SthreeEThreeInputChangev, and the cost of a run by \SthreeEThreeCostv: less than the tokens, because most of the
input re-read at each turn is billed at the cached price. H8 fails on this case because of that one run: the
strategy's median, \SthreeEThreeTurnsSv, is not below \SthreeEThreeTurnsRminv.

\paragraph{Billing support, where it already grouped its calls.} Little was left: runs without the strategy took
\SthreeEThreeTurnsRminiv{} to \SthreeEThreeTurnsRmaxiv{} turns, those with it always \SthreeEThreeTurnsSiv, which
is the fewest this task needs (see the post-hoc analysis below).

\paragraph{Nothing was lost for it, and the prompts were not overridden.} H10 holds: the gates passed in
\SthreeHTenGatesiv{} and \SthreeHTenGatesv{} runs of both arms of billing support and CI triage, every requirement of
the prompts covered. H11 holds: the prescribed values were respected in \SthreeHElevenPassed{} of
\SthreeHElevenChecks{} checks, in the \SthreeHElevenRuns{} strategy runs with recorded outputs. Those two cases point
the same way as CI triage (median turns from \SthreeEThreeTurnsRi{} to \SthreeEThreeTurnsSi{} in the editorial
plan, from \SthreeEThreeTurnsRii{} to \SthreeEThreeTurnsSii{} in the content-gap analysis), on
\SthreeEThreeRunsRi{} and \SthreeEThreeRunsRii{} runs per arm and with recorded outputs: the model's call matched a
recorded call exactly in \SthreeExactMini{} to \SthreeExactMaxi{} of the calls of a run in the editorial plan and in
\SthreeExactMinii{} to \SthreeExactMaxii{} in the content-gap analysis, and got the nearest recorded output otherwise.
They say how calls were organised, not whether the task was done well.

\paragraph{Post hoc: calls that could have been made earlier.} This analysis is not in the protocol; it was added
after the runs, to tell a sequence the task imposes from one the model chose. A call whose arguments name something
an earlier result returned (a run, a job, a customer, an invoice) can only follow that result; any other could have
been made at once. The longest chain of such calls gives the fewest turns a run needed. In CI triage
\SthreeDependingv{} of the calls depended on an earlier result, but the chain was short (the runs, then their jobs,
then the logs): \SthreeFewestv{} turns were enough in every run. Without the strategy \SthreeEarlierRv{} of the calls
were made later than they could have been, with it \SthreeEarlierSv; in billing support \SthreeEarlierRiv{} and
\SthreeEarlierSiv, with \SthreeFewestiv{} turns needed. The rule reads identifiers only and cannot know that a model
wanted to read one log before asking for the next: it gives a lower bound on the turns, not what a careful agent
should do.

\subsection{Offline analyses, not in a protocol}
\label{sec:s3offline}

Three recalculations on runs already recorded, made after E3 with no new run. None was in a protocol; they are
reported as such.

\paragraph{Projections per mode of use.} Would one projection per value of a low-cardinality argument (a result
type, a time range, a sort order) do better than one projection per tool? On the \SthreeModesRuns{} recorded runs
in which the model saw whole outputs, each held out in turn, \SthreeModesCandidates{} candidate arguments on
\SthreeModesTools{} tools were evaluated, identifiers left out. The rule, written before the numbers: a benefit is at
least \SthreeModesThreshold{} more points of the tool's tokens removed without more calls losing a field the run
used. A benefit was found for \SthreeModesBenefit{} candidate. The most any candidate added was
\SthreeModesMaxPoints{} points, and \SthreeModesUnderOne{} added less than \SthreeModesSmall; \SthreeModesLoseMore{}
lost a used field in more calls than the projection per tool, because a profile learns from fewer calls;
\SthreeModesThin{} had a profile with fewer than \SthreeModesMinProfile{} calls, and \SthreeModesFallback{} held-out
calls fell on a profile too thin to have its own projection. The specialisation was not implemented.

\paragraph{Rows.} The library removes fields, never rows. In the \SthreeRowsRuns{} recorded runs in which the model
saw whole results, on the search-data API the elements of a list that the run reused afterwards were exactly the
first ones in \SthreeRowsTopSearch{} of the calls with any reuse, and keeping only the first \SthreeRowsCut{} would
have removed a reused element in \SthreeRowsCutFiveSearch{} of them; on the other APIs, whose lists are short, the
figures are \SthreeRowsTopOther{} and \SthreeRowsCutFiveOther. On the search-data API a cut by position would remove
elements the agent reuses, and a removed row, unlike a removed field, cannot be told from one that does not exist.
No cut of rows is implemented; the telemetry report shows how many list elements are reused, and from where.

\paragraph{Tool definitions.} The model receives the definition of every listed tool at every turn. Counted on the
schemas as sent in the \SthreeDefRuns{} recorded runs of E1, E2 and E3, in the reference arms the definitions were
\SthreeDefInputMin{} to \SthreeDefInputMax{} of the input tokens (median per case and model), but, sitting in the
cached head of the prompt, \SthreeDefCostMin{} to \SthreeDefCostMax{} of a run's cost at the providers' list prices;
had nothing been cached, up to \SthreeDefUncachedMax. A counterfactual, with no run made, asks what a whole server
would weigh: the search-data provider's MCP server lists \SthreeDefServerTools{} tools, of which the editorial-plan
agent was given \SthreeDefCaseTools{} and called \SthreeDefUsedTools. Taking its \SthreeDefScenarioRuns{} recorded
R runs as they happened and changing only the definitions, listing every tool would send \SthreeDefAllTokens{}
tokens per turn instead of \SthreeDefCaseTokens; listing the tools called and putting the other
\SthreeDefHiddenTools{} behind a search tool and a call tool, with no search ever made (the most hiding could gain),
\SthreeDefDeferredTokens, \SthreeDefSavedShare{} fewer than the whole server, and the first turn would go from
\SthreeDefFirstTurnAll{} to \SthreeDefFirstTurnHidden{} tokens. With a cache, depending on what is assumed cached,
the cost of a run would fall by \SthreeDefHideGainMin{} to \SthreeDefHideGainMax, and \SthreeDefExtraTurnsMin{} to
\SthreeDefExtraTurnsMax{} extra turns (medians) would cancel the gain; the scenario also assumes that the model
works with the whole server as it did with its own tools. The library therefore measures the definitions and does
not hide tools: its report lists the tools never called, with their weight, and suggests disabling them in the host
or using native search where it exists. Some hosts already defer definitions: Claude Code loads MCP tools on demand
by default \citep{anthropic2026mcp}; the OpenAI Responses API offers tool search when the request asks for it
\citep{openai2026toolsearch}; Cursor does so according to a blog post, which we did not verify in its documentation
\citep{cursor2026dynamic}. A second search on top of theirs would add turns, and a hidden tool called through a
generic call tool would bypass the per-tool permissions that hosts grant.

\subsection{Hosts}
\label{sec:s3hosts}

Between 11 and 13 November 2026 the proxy was put in front of servers in several hosts, in most checks with a log that
records each request of the host's client and none of its content; the repository's page on hosts states what was verified,
when and how, and everything else as not verified. In the OpenAI Playground, Claude Code (on Linux and on Windows)
and VS Code (Copilot Chat in agent mode), the tools passed and the model called \texttt{tracetrim\_done} at the end
of its run; in the Playground it did so guided by the tool's description alone. \texttt{expand} was verified in the
Playground only: in the other hosts nothing had yet been trimmed. MCP Inspector, which has no model, showed that the
skill is served as the MCP Skills extension asks. None of the hosts with a model loaded the skill as a skill: no
client declared the extension. The Playground's and VS Code's clients asked for neither the skill nor resources; in Claude Code the model
read the skill through the host's tool for MCP resources before its first call, as a resource. Claude Desktop and
Cursor were not verified. These checks show that the path works in those hosts, not how often or how well.
